% !TeX root = main.tex
\section{未来研究方向}\label{sec:future}

从前文的讨论可以看出，给定偏好下的多目标序贯决策，需要把长期准则、信息表示、策略求解与最终评价贯通起来。近期研究已在非线性效用的近似规划、跨偏好策略改进以及非线性离线优化等方面取得进展\cite{peng2025esr,alegre2026,kim2025fairdice,kim2026aetdice}。接下来的问题是：这些结果在什么准则和信息条件下成立，近似表示与有限数据带来的误差如何影响原偏好下的性能，求得的策略在部署条件下能否保持相应保证。本章据此提出四个研究方向，分别针对信息表示、跨偏好复用、固定数据条件和部署条件中尚未解决的问题。

\subsection{复杂长期准则的充分表示与近似误差控制}\label{sec:future-representation}

复杂偏好准则难以求解，原因既在于评价需要的信息多，也在于保留这些信息的计算代价高。第~\ref{sec:structure-recursion}~节说明，非线性 ESR 可以在包含累计回报和时间的扩展状态上形成递推。Peng 等进一步在有限状态动作、有界确定性向量奖励、固定目标维数和效用一致连续等条件下，建立了考虑累计结果的 Bellman 最优性关系，并给出近似规划及结合模型学习的求解方法\cite{peng2025esr}。这一结果说明非线性效用并不排斥动态规划；其效率则依赖目标维数、效用连续性、时域或折扣截断等条件，累计回报网格的复杂度对目标数呈指数依赖（第~\ref{sec:solving-model}~节）。

问题由此从“递推是否存在”推进到“递推需要保留多少信息”。对给定效用而言，完整的累计向量或完整的回报分布未必是最经济的表示；但只保留均值、逐目标的边缘分布或几个低阶统计量，又可能改变原来的策略排序。例如，ESR 中多个目标共同达标的效用需要目标之间的联合信息，非线性 SER 的策略比较也无法由逐状态施加效用的局部选择代替\cite{hayes2022,ropke2023}。因此，表示压缩要围绕具体准则定义充分性，分别检查它是否保留策略评价所需的信息、是否支持最优性递推、是否允许原问题需要的策略形式。

一个具体的研究问题是：针对给定的效用类别和路径要求，能否构造规模较小、决策误差可控的充分表示。可以从累计量的等价类、与效用相关的统计量或任务专用的记忆入手，分析被合并的历史在后续可行动作下是否给出相同的比较结果；只能近似合并时，再把信息损失转化为原效用下的策略性能误差。对于连续或敏感度可控的效用，可以按效用对不同回报区域的敏感程度分配离散精度，而不是对所有目标和累计水平使用同样的网格。阈值事件和严格优先关系需要单独处理，因为阈值附近很小的表示误差就可能改变可接受性或比较次序。

这一方向还要把表示误差与模型误差、价值估计误差和数值求解误差区分开。有限时域偏好关系理论已经给出递归最优性的结构条件\cite{carr2024}；把这些条件落实到连续状态或高维近似表示时，需要判断近似后的子问题是否仍保持原有的比较关系。研究进展可以用原准则下的性能误差随记忆规模、目标数和时域的变化来衡量：先在能精确求解的小规模问题中验证误差界，再考察它在函数逼近中的有效性。训练损失下降或解集超体积提高，都不足以证明表示压缩保持了给定偏好。

\subsection{跨偏好复用的表达范围与性能保证}\label{sec:future-reuse}

跨偏好复用的核心问题，是共享的策略和价值信息能为哪些给定偏好提供可靠的解。Alegre 等围绕线性效用研究广义策略改进，根据候选策略的覆盖情况给出相对于最优解的效用损失界，并在置信下界有效的条件下，控制不确定性感知策略改进相对于库中最好策略的损失\cite{alegre2026}。线性偏好下的策略覆盖、样本复用和近似误差因此已有可分析的方法基础，而这些保证所依赖的线性回报组合、候选策略和价值估计条件，也划定了可以直接复用的准则范围。

非线性偏好下也已有实质进展。Kubo 等的近期公开稿针对有限状态动作、折扣准则和平稳随机策略下的平滑 Tchebycheff 效用，在偏好参数位于规定内部区域等条件下，证明了最优回报向量的唯一性及其关于偏好的 Lipschitz 连续性，并在占用测度上构造了具有收敛率的迭代方法\cite{kubo2026}。这为特定非线性 SER 下的偏好扫描和条件策略学习提供了理论依据。需要注意，最优回报向量连续不等于最优策略映射连续，精确的表格迭代保证也不直接适用于容量有限的深度条件网络。

Kim 等对严格正有界的逐阶段奖励、固定有限时域以及给定效用的最优求解器，在广义均值族内分别建立了 SER 与 ESR 的策略组合近似保证\cite{kim2025portfolio}。后续问题因此可以更具体：效用类别超出该族、奖励允许零值或负值，或者底层求解器只有有限数据下的近似保证时，覆盖误差如何变化；精确 oracle 下的候选规模界，也要与获得每个候选的实际成本和近似误差一起分析。

后续研究需要明确复用对象的表达范围。对于非线性 SER，线性偏好所需的凸覆盖信息不一定支持受限策略类别下的全部非线性比较，允许策略混合时可达集合还会改变；对于 ESR，还要保留评价所依赖的回报分布信息。Röpke 等针对规定的多变量风险厌恶效用类别建立了分布支配和候选保留机制\cite{ropke2023}，其覆盖性质限于该效用类别。策略库、条件网络和共享价值表示，都应说明自己能表达哪些策略，以及偏好类别扩大后需要补充哪些信息。

偏好覆盖与计算资源之间的定量关系值得进一步研究，可以以原准则下的效用损失为依据，分析偏好参数变化、候选策略覆盖误差、价值估计误差和条件策略逼近误差的共同作用；在能够获得最优值或可靠参考界的条件下，考察哪些偏好区域可以直接复用、哪些区域需要增加候选或专门求解。对效用敏感度高、约束活跃集变化或策略选择发生切换的区域，可以有针对性地加密偏好采样、分配局部容量，而不是在权重空间中追求均匀覆盖。这些机制的效果，应同时用平均偏好和较差偏好区域的效用损失、新偏好的追加求解成本以及部署计算量来检验。

\subsection{固定数据下的偏好相关覆盖与可信评价}\label{sec:future-offline}

有限历史数据下的多目标优化已经超出了线性回报重加权。FairDICE 直接优化各目标期望回报上的严格凹福利函数，通过占用分布正则化和分布校正估计构造离线策略，其理论对应带正则化的非线性 SER 问题\cite{kim2025fairdice}。近期的 AETDICE 公开稿通过有限时域的累计回报与时间扩展，在外层聚合为凹函数等条件下，把 SER、ESR 及一定形式的组合准则纳入同一离线优化框架\cite{kim2026aetdice}。研究问题因此已经从“非线性多目标准则能否用固定数据求解”，转向“所需信息和覆盖条件能在有限日志中得到多大程度的支持”。

累计回报扩展让这个问题更加具体。AETDICE 的理论构造要求扩展状态动作上的数据覆盖，并指出累计回报维度增加会加剧覆盖困难\cite{kim2026aetdice}。同一状态和动作在数据中频繁出现，不代表它们在某个累计结果或剩余时域下也得到了充分支持；不同偏好又可能把策略带到不同的累计回报区域。因此，原状态动作覆盖、行为偏好覆盖和扩展信息覆盖需要分别分析。D4MORL 对不同偏好分布的比较实验为此提供了基准基础\cite{zhu2023}，但改变偏好标签或重用同一条轨迹，都不会增加未观测动作的真实后果信息。

可以进一步研究由数据和准则共同决定的可靠求解范围：在给定偏好下，哪些策略结果能从当前日志中识别，误差控制需要什么样的覆盖条件，数据不足时应返回什么样的保守解或“无法判定”。对于可平滑评价的效用，可以分析各目标的估计误差及其相关性如何传播为效用损失；对于 ESR、联合达标事件和尾部风险，则需要针对联合累计结果或相应的轨迹统计量建立估计。保守机制应直接约束原准则下的估计风险：逐目标分别做悲观估计，并不等价于对整体效用或联合安全要求做悲观估计。同时还应区分行为正则化带来的优化偏差、有限样本误差，以及从占用分布中提取可执行策略的误差。

可靠的评价也关系到离线模型选择。正则化强度、策略容量和偏好覆盖机制都会影响最后选出的策略；如果这些选择只能依靠新的环境执行结果来确定，整个决策流程使用的信息就已经超出了固定日志。均值和标量分布的离线评价方法为估计长期结果提供了基础\cite{jiang2016,chandak2021}，条件联合回报分布估计进一步保留了目标之间的相依关系，但其理论依赖全支持、遍历性和函数类等条件\cite{qi2025distributional}（第~\ref{sec:eval-offline}~节）。后续需要在这些结果的基础上，控制同时比较多个策略、多个偏好并据此选择模型所带来的误差，并弄清覆盖薄弱时哪些结论仍然可信。独立验证日志或交叉拟合是可行的路径，但仍需写明目标策略的支持、模型误差和所需的轨迹信息。

这一方向的实验，应独立改变状态动作支持、累计结果支持、偏好范围和目标数，并报告原效用误差、策略排序错误和约束误判，而不只是改变专家轨迹的比例。提出置信保证时，还应检验置信区间的实际覆盖率、覆盖薄弱区域的表现以及无法可靠判定的比例。这样才能区分算法改进、数据条件改善和额外环境信息各自带来的收益，确定固定历史数据实际能支持的决策范围。

\subsection{部署信息与模型不确定性下的闭环保证}\label{sec:future-deployment}

求解时的信息充分性需要延续到实际执行。Vamplew 和 Foale 的近期研究札记指出，依赖累计回报的扩展状态策略在停止学习后，仍可能需要持续获取实际回报或能够支持相应决策的代理信号；其 ESR 反例显示，用不可观测随机奖励的均值代替实际累计结果，可能改变后续的动作选择和原效用下的表现\cite{vamplew2026deployment}。这揭示了状态扩展与部署之间的具体接口：训练日志中的奖励变量在部署时不一定可观测，学习阶段构造了充分状态，不代表执行阶段拥有同样的信息。

后续可以把策略表示从训练时可用的变量转向部署时可获得的信息，研究可观测的累计量、延迟反馈和隐藏回报的联合信念是否足以支持原准则，并分析代理信号的误差如何影响效用和约束判断。若回报能由可观测的状态、动作和转移精确恢复，累计量仍可正常更新；若实际随机回报无法由部署观测确定，即使准确知道它的条件均值或分布，也可能无法执行完整信息下的最优策略。因此，策略应定义在部署可用的信息集或相应的联合信念上，分别分析回报不可观测造成的信息损失和代理模型的估计误差，并在训练与评价中保持这一信息条件。部署验证可以通过回报缺失、反馈延迟和代理误差的受控实验，把策略本身的求解误差与执行信息变化造成的性能损失区分开。

环境模型的不确定性还会改变长期结果及其可接受性。这里要区分目标权衡、随机风险和模型不确定性三个层次。以给定偏好 $p$ 和预先规定的模型集合 $\mathcal M$ 为例，可以考察
\begin{equation}
J_p^{\mathrm{rob}}(\pi)
=\inf_{M\in\mathcal M}J_p^M(\pi),
\end{equation}
其中 $J_p^M$ 沿用模型 $M$ 下原有的 SER、ESR 或其他长期评价定义。不同目标的最坏结果可能出现在不同的模型下，因此“逐目标取最坏值再做效用评价”与“同一模型下的整体最坏效用”是两个需要分别界定的量。Cialenco 和 Kováčová 针对与向量结构相容的序关系构造了最坏期望损失及集合值函数，并在相应的矩形性等条件下建立了较强的递推关系\cite{cialenco2026robust}。这类最坏向量与上式按整体效用取最坏的对象未必相同，其结果适用于他们定义的向量序关系，用于非线性 SER 或 ESR 的鲁棒保证还需要专门分析；不确定性在阶段间和目标间的耦合方式也要与原准则一起考虑。完整结果上的风险评价与嵌套条件风险，同样要按各自的时间语义分别处理\cite{ruszczynski2010}。

控制研究已经为部分闭环要求提供了具体保证。Eichfelder 等把多目标模型预测控制中的耗散性与兼容终端代价条件放宽到只对一个目标成立，并在相应的选择规则和附加条件下分析了闭环稳定性与多目标性能\cite{eichfelder2023}；Wu 等针对已知的名义模型和有界扰动，通过约束收缩、终端设计和辅助稳定性约束建立了递归可行性与输入到状态稳定性\cite{wu2024}。因此，这一方向的研究问题不是笼统地“增加鲁棒性”或“走向真实应用”，而是：给定偏好下的原始效用、物理可行性和闭环保证如何同时成立，为了获得稳定性而限制候选解时，要付出多少偏好性能代价。

一种可行的做法是，把偏好求解器与状态估计、模型更新和安全动作修正组成的完整执行策略作为分析对象，明确扰动范围、估计信息和实时计算条件，再研究效用性能界与可行性保证同时成立的条件。原准则不能直接写成预测时域内可加代价时，还要检查终端表示和滚动重优化是否保留了累计结果和跨期取舍。验证时应同时报告原准则下的效用、约束违反、模型失配范围和决策延迟，并比较安全修正前后的实际表现。这一方向要回答的是：理论上可以求解的偏好准则，经过信息估计、近似规划和实际执行之后，还能保留哪些性质。
